% Tipo de documento. En este caso es un art�culo, para folios A4, tama�o de la fuente 11pt y con p�gina separada para el t�tulo
\documentclass[a4paper,11pt,titlepage]{article}

% Carga de paquetes necesarios. OrdenesArticle es un paquete personalizado
\usepackage[spanish]{babel} 
\usepackage{eurosym} %para poder poner el s�mbolo del euro con \euro
\RequirePackage[T1]{fontenc}
\RequirePackage[ansinew]{inputenx} 
\usepackage[spanish,cap,cont,title,fancy]{OrdenesArticle}
\usepackage{lmodern}
\usepackage{graphicx}
\usepackage{hyperref}
\usepackage[usenames,dvipsnames]{color}
\usepackage{color}
\usepackage{ifthen}
\usepackage{longtable}
\usepackage{listings}
\usepackage{enumerate}
\hypersetup{bookmarksopen,bookmarksopenlevel=3,linktocpage,colorlinks,urlcolor=blue,citecolor=blue,
						linkcolor=blue,filecolor=blue,pdfnewwindow,
						pdftitle={FAQ Manager: un gestor de preguntas frecuentes},
						pdfauthor={Juan Andrada Romero, Jose Domingo L�pez L�pez},
						pdfsubject={Multimedia, Hipermedia y Realidad Virtual}}


% Macro para definir una lista personalizada 
\newenvironment{milista}%
{\begin{list}{\textbullet}%
{\settowidth{\labelwidth}{\textbullet} \setlength{\leftmargin}{\dimexpr\labelsep+\labelwidth+5pt}
\setlength{\itemsep}{\dimexpr 0.5ex plus 0.25ex minus 0.25ex}
\setlength{\parsep}{\itemsep}
\setlength{\partopsep}{\itemsep}
\addtolength{\topsep}{-7.5pt}
}}%
{\end{list}}

% Macro para insertar una imagen
%       Uso: \imagen{nombreFichero}{Factor escala}{Caption (leyenda)}{Label (identificador para referenciarla)}
% -------------------------------------------------------------------------------------------------------------
\def\imagen#1#2#3#4{
 \begin{figure}[h]
 \begin{center}
   \scalebox{#2}{\includegraphics{#1}}
 \caption {#3}
 \label{#4}
 \end{center}
 \end{figure}
}

% Definicion del "listings" para el lenguaje Java
\lstdefinelanguage{Java}
{
 morecomment = [l]{//}, 
 morecomment = [l]{///},
 morecomment = [s]{/*}{*/},
 morestring = [b]", 
 sensitive = true,
 morekeywords = {package, static, while, switch, break, line, void, String, Object, int, Integer, instanceof, else, if, for, private, return, new, public, class, import, int, boolean, true, false, extends, final, super, protected, abstract, this, do, float, double, null, try, catch, implements}
}

% Configuraci�n del c�digo incluido para el lenguaje Java
\lstset{
  language=Java,
  basicstyle=\footnotesize,
  backgroundcolor=\color{white},
  showspaces=false,
  showstringspaces=false,
  showtabs=false,
  frame=single,
  tabsize=2,
  captionpos=b,
  breaklines=true,
  breakatwhitespace=false,
  escapeinside={\%},
  keywordstyle = \color [rgb]{0,0,1},
  commentstyle = \color [rgb]{0.133,0.545,0.133},
  stringstyle = \color [rgb]{0.627,0.126,0.941}
}


\begin{document}

% En las p�ginas de portada e �ndices, no hay encabezado ni pie de p�gina
\pagestyle{empty} 

% Se incluye la portada
\begin{titlepage}
	\begin{center}
  	{\LARGE UNIVERSIDAD DE CASTILLA-LA MANCHA} \\
  	\bigskip
  	{\Large ESCUELA SUPERIOR DE INFORM�TICA} \\
  	\vspace{20mm}
  	\includegraphics[scale=0.45, keepaspectratio]{./images/esi_bw.png} \\
  	\vspace{20mm}
  	{\Huge \textbf{Almacenamiento y Recuperaci�n de la Informaci�n}}\\
  	\vspace{10mm}
  	{\LARGE \textsc{\textbf{M�dulo de importaci�n e indexaci�n de documentos}}}\\
  	\vspace{30mm}
  	\today
  	\vspace{40mm}
  	\begin{flushleft}
  		{\large Juan Andrada Romero}\\
  		\vspace{1mm}
  		{\large Jose Domingo L�pez L�pez}\\		
  	\end{flushleft}
	\end{center}
\end{titlepage}
\clearpage

% Texto del reverso de la portada
%%\mbox{}
%%\vspace{18cm}
%%\begin{small}
% Se ajusta la separaci�n entre p�rrafos
%%\parskip=10pt 

%%\copyright~ 2008/2009 Juan Andrada Romero. Universidad de Castilla La Mancha, Escuela Superior de Inform�tica de Ciudad Real.

%%Se permite la modificaci�n, copia y distribuci�n de este documento, seg�n la licencia de documentaci�n GNU (\url{http://www.gnu.org}).

%%Este documento fue compuesto con \LaTeX{}. Im�genes generadas con Power Point y Gimp.
%%\end{small}

%%\newpage

% En las p�ginas de �ndices y prefacio, se utiliza numeraci�n romana
\pagenumbering{Roman}

% Se crea el �ndice
\tableofcontents
% Se pasa p�gina y se a�ade el �ndice de figuras al �ndice principal
%\clearpage\phantomsection
%\addcontentsline{toc}{section}{\listfigurename}
% Se crea el �ndice de figuras
%\listoffigures
% Si se quiere crear un �ndice de tablas se pondr�a: \listoftables

\newpage

% Se ajusta la separaci�n entre p�rrafos
\parskip=10pt

% Se a�ade el prefacio al �ndice
%\clearpage\phantomsection
%\addcontentsline{toc}{section}{Prefacio}

% Comienza el contenido del documento. Se utilizan n�meros ar�bigos y el encabezado y pie de p�gina personalizado
\pagenumbering{arabic}
\pagestyle{fancy}

\section{Participantes}

La pr�ctica ha sido desarrollada por los siguientes alumnos:
\begin{milista}
	\item Juan Andrada Romero (grado de participaci�n 50\%)
	\item Jose Domingo L�pez L�pez (grado de participaci�n 50\%)
\end{milista}
%\clearpage
\section{Introducci�n}

Cuando un usuario visita una p�gina Web, puede plantearse infinidad de preguntas acerca de sus caracter�sticas, pol�tica de privacidad, t�rminos de uso, problemas con un producto, etc. En ocasiones, varios usuarios formulan las mismas preguntas diario y las ponen en conocimiento del personal administrativo de la p�gina Web, siendo necesario responder todas y cada una de ellas de un modo individual y personalizado. En este momento el propietario del sitio debe plantearse en la posibilidad de a�adir una secci�n de \textit{preguntas frecuentes}, tambi�n conocidas como \textbf{FAQ} (\textit{Frequently Asked Questions}).

\textbf{FAQ Manager} es un \textit{gestor de preguntas frecuentes} construido en \textbf{Flash}, que trata de ofrecer una soluci�n al problema planteado de una forma �gil, din�mica y con proceso de instalaci�n sencillo. Permite al personal administrativo crear un conjunto de categor�as bajo las cuales podr�n incorporar una serie de preguntas junto a sus respuestas (cada pregunta debe pertenecer obligatoriamente a una categor�a). Una vez configuradas las categor�as, las preguntas y las respuestas, la p�gina principal del sistema permitir� a los usuarios navegar entre las distintas categor�as para leer las preguntas que le corresponda.
\clearpage
\section{Descripci�n de la base de datos}

La persistencia de \textbf{FAQ Manager} consiste en dos tablas que deber�n ser creadas en la base de datos del sistema:
\begin{milista}
	\item \textit{categoria\_faq} es la tabla que almacenar� las categor�as en las que se clasificar�n las \textit{preguntas frecuentes}. Contiene dos atributos: un identificador autonum�rico y el nombre de la categor�a.
	
	\item \textit{faq} es la tabla que almacenar� las \textit{preguntas frecuentes} y sus respuestas correspondientes. Consta de cuatro campos: un identificador autonum�rico, la pregunta, la respuesta y el identificador de la categor�a a la que pertenece la pregunta.
\end{milista}

En la distribuci�n de \textit{FAQ Manager} puede encontrar un fichero denominado \texttt{install.sql}. Este fichero es un \textit{script} que contiene las sentencias necesarias para crear una base de datos llamada \textit{mhrv} y las dos tablas mencionadas anteriormente.

%\clearpage
\section{Requisitos del sistema}

Para poder utilizar \textbf{FAQ Manager} ser� necesario disponer de un servidor de p�ginas Web (por ejemplo, Apache HTTP Server \cite{apache-homepage}) que pueda procesar c�digo PHP \cite{php-homepage}. Adem�s, ser� necesario un sistema gestor de bases de datos como MySQL \cite{mysql-homepage}. No obstante, para simplificar el proceso de instaci�n y configuraci�n de todo este software, se recomienda el uso de WAMP \cite{wamp-homepage}.

Una vez instalado y configurado el software necesario, deber� editar el fichero \texttt{config.php} que se incluye en la distribuci�n de \textit{FAQ Manager}. En este fichero, de estructura similiar a la de Figura \ref{fig:config}, deber� cumplimentar debidamente el dominio sobre el que se aloja el sistema gestor de bases de datos, el nombre de la base de datos, el nombre de un usuario de la base de datos con los permisos b�sicos para realizar operaciones \textit{CRUD} (Create, Read, Update, Delete) y su contrase�a.

\imagen{./imagenes/config.png}{0.7}{Fichero \texttt{config.php}}{fig:config}

\clearpage
\section{Manual de usuario}

Es importante destacar que la aplicaci�n se ha desarrollado desde la perspectiva del administrador, ya que no se ha implementado ning�n m�dulo de \textit{autenticaci�n} para gestionar una pol�tica de acceso en la que los visitantes no \textit{autenticados} no puedan modificar el estado de la aplicaci�n.

\subsection{Entrada en la aplicaci�n y descripci�n general}

Para entrar en la aplicaci�n es necesario visitar la direcci�n \cite{faq-manager} desde un navegador Web. Inicialmente, las tablas de la base de datos que utiliza la aplicaci�n est�n vac�as, por lo que no contienen informaci�n para mostrar (ver Figura \ref{fig:inicio-vacio}) y deber� ser configurada por el administrador del sistema.

\imagen{./imagenes/1inicio.png}{0.65}{Pantalla inicial vac�a}{fig:inicio-vacio}

En la esquina superior izquierda de la pantalla principal (ver Figura \ref{fig:inicio-vacio}) se observa un bot�n denominado \textit{Configuraci�n}. Este bot�n nos dirigir� a una nueva pantalla correspondiente al \textit{Panel de Control} (ver Figura \ref{fig:panel}) desde el que se podr� a�adir nuevas categor�as, a�adir nuevas preguntas (con sus respuestas), y modificar y eliminar preguntas ya existentes en la base de datos.

\imagen{./imagenes/2panelcontrol.png}{0.65}{Panel de control}{fig:panel}


\subsection{A�adir nuevas categor�as y FAQ}

Para a�adir una nueva categor�a es necesario situarse en el \textit{Panel de Control} y escribir el nombre de la categor�a que se desea crear en la caja de texto reservada para este prop�sito. A continuaci�n se pulsa el bot�n \texttt{A�adir} y aparecer� un mensaje verde, indicando si la categor�a se ha creado correctamente, o rojo, si se ha producido alg�n error.

Inmediatamente despu�s de crear una nueva categor�a (ver Figura \ref{fig:catadd}), ya es susceptible de a�adir \textit{preguntas frecuentes} bajo la misma, ya que las cajas de selecci�n de las operaciones \textit{A�adir FAQ} y \textit{Modificar o eliminar FAQ} se refrescan autom�ticamente, no siendo necesario refrescar la p�gina.

\imagen{./imagenes/3catadd.png}{0.65}{Categor�a a�adida y p�gina refrescada autom�ticamente}{fig:catadd}

Del mismo modo, se pueden crear nuevas \textit{preguntas frecuentes} utilizando la regi�n dedicada para este prop�sito (la segunda regi�n). Los tres campos propuestos (pregunta, categor�a y respuesta) son obligatorios, por lo que no se podr� crear ninguna pregunta hasta que haya sido dada de alta, al menos, una categor�a. Se desea destacar la posibilidad de escribir texto \textbf{HTML}, ofreciendo as� un mayor poder expresivo al permitir m�s flexibilidad a la hora de dar formato al texto (uso de negrita, cursiva y subrayado), inclusi�n de enlaces e inclusi�n de im�genes (ver Figura \ref{fig:faqaddimg}) y listas no numeradas.

\imagen{./imagenes/4faqaddimg.png}{0.65}{A�adir una pregunta con c�digo HTML}{fig:faqaddimg}


\subsection{Modificar o eliminar FAQ}

Para modificar o eliminar una FAQ, primero deber� seleccionar la pregunta deseada en la lista desplegable destinada para dicho prop�sito (la tercera regi�n). Una vez que seleccione una pregunta, el cuadro se refrescar� autom�ticamente recuperando del servidor la categor�a y la respuesta actual.

Observe que aparecen dos botones: \texttt{eliminar} y \texttt{modificar}. En este momento el usuario puede \textbf{eliminar} la pregunta (y su respuesta) de la base de datos o puede \textbf{modificar} la informaci�n de dicha pregunta seleccionando otra categor�a distinta (para cambiar la categor�a bajo la que se alojar� la pregunta) y/o modificando el contenido de la respuesta para rectificar alg�n error y/o a�adir o eliminar informaci�n. Una vez m�s, se permite la escritura utilizando c�digo HTML para enriquecer el formato del texto.


\subsection{Consulta de FAQ}

El proceso de consulta y lectura de FAQ es muy sencillo. El usuario deber� situarse en la p�gina principal del sistema. En ella se observan dos paneles (ver Figura \ref{fig:fin}):
\begin{milista}
	\item En el panel de la \textbf{izquierda} aparecer�n las distintas categor�as que han sido dadas de alta en el sistema.
	\item En el panel de la \textbf{derecha}, inicialmente vac�o, ir�n apareciendo preguntas y respuestas conforme vayamos seleccionando las distintas categor�as del panel izquierdo.
\end{milista}

\imagen{./imagenes/5inicio.png}{0.65}{Sistema configurado y en uso}{fig:fin}
\clearpage

% Secci�n de ap�ndices
% Se a�ade el titulo de "Ap�ndices" al �ndice
%\addcontentsline{toc}{section}{Ap�ndices}

%\input{Apendices.tex}

% A�adimos la bibliografia al �ndice
\clearpage\phantomsection
\addcontentsline{toc}{section}{Referencias}

% Bibliograf�a. En este caso se usa BibTeX
\bibliographystyle{plain}
\pagestyle{plain} 
\bibliography{Bibliografia}

\end{document}
